% MT003 / explanation — canonical editable content.
本题属于已保留的历史拓展。高斯过程（Gaussian process，GP）直接给未知函数赋予概率模型：取任意有限个输入，对应的函数值共同服从一个多元高斯分布。均值函数描述大致水平，核函数描述不同位置的函数值怎样一起变化。



\begin{itemize}

\item \textbf{A不选。} RBF是常用核，不是GP定义。GP也可以采用其他合法核；不能把“贝叶斯线性回归换成RBF”当作全部定义。

\item \textbf{B按本课标准模型选。} 高斯函数先验配合高斯观测噪声，观测后的函数分布仍是高斯，可以通过矩阵公式求均值和协方差，这就是常说的解析后验。原题的“只有”不宜推广为所有可能特殊模型的数学定理。

\item \textbf{C按RBF设定选。} RBF核给近输入更大的协方差，因此观察附近一个函数值会更多地影响当前点的预测。但“近就更相关”由核决定，并非每一种GP核都必须如此。

\item \textbf{D不选。} 标准精确GP的大数据难点是要处理\(N\times N\)协方差矩阵；稠密计算通常需\(O(N^3)\)时间及\(O(N^2)\)存储，不是因为参数少所以表达能力不够。

\item \textbf{E不选。} 输入可由实验者固定，也可按时间顺序采集；GP并不要求所有输入位置独立同分布。标准模型常假设观测噪声独立，这与输入独立是两回事。

\end{itemize}


\textbf{出处：} 2020B期中Quiz Q3，原卷第2页（5分）


\textbf{答案依据：} 2020B期中 Q3 随卷／配套参考答案；条件与解释经复核。
